% =============================================================================
% writeup-template.tex
%
% ACCT 932 -- Beaver (1968) replication write-up, LaTeX version.
%
% HOW THIS WORKS
%   You do not paste numbers into this document. The pipeline writes
%   tables to OUTPUT_DIR as .tex files and figures as .pdf, and this
%   document pulls them in with \input{} and \includegraphics{}. Re-run
%   the pipeline, recompile, and every number updates itself. That is the
%   entire point -- a transcription error becomes impossible.
%
% TO COMPILE
%   Locally:   pdflatex writeup-template  ->  bibtex writeup-template
%              ->  pdflatex twice more
%   Overleaf:  upload this file, references.bib, and the contents of your
%              output/ folder. Overleaf's free tier does not sync with
%              GitHub, so you upload output/ by hand after each run.
%
% NOTE ON PATHS
%   \graphicspath below assumes the default layout, with this file in
%   writeup/ and OUTPUT_DIR set to output/. If you changed OUTPUT_DIR in
%   .env, change the paths here to match.
% =============================================================================

\documentclass[12pt]{article}

\usepackage[margin=1in]{geometry}
\usepackage{graphicx}
\usepackage{booktabs}
\usepackage{caption}
\usepackage{float}
\usepackage{amsmath}
\usepackage[colorlinks=true, allcolors=blue]{hyperref}
\usepackage{natbib}
% plainnat ships with natbib itself, so this compiles on a bare TinyTeX
% install as well as on Overleaf. If you prefer a journal's house style
% (chicago, apalike, aer), swap it here -- but on a local TinyTeX you may
% first need: tinytex::tlmgr_install("chicago")
\bibliographystyle{plainnat}

\usepackage{xcolor}

% -----------------------------------------------------------------------------
% WRITING PROMPTS
%
% \guidance{} holds the instructions telling you what to write in a
% section. It renders as grey italic text, so a prompt you have not yet
% answered is obvious on the page. Delete each one as you replace it with
% your own writing.
%
% These are body text rather than % comments because this file is the
% single source for the Word version too: writeup/make-writeup-docx.R
% converts this document, and a LaTeX comment would not survive that
% conversion. Anything a student needs to READ belongs in \guidance{};
% anything only a LaTeX user needs belongs in a % comment.
% -----------------------------------------------------------------------------
\newcommand{\guidance}[1]{{\itshape\color{gray}#1\par\medskip}}

% -----------------------------------------------------------------------------
% REQUIRED PREAMBLE FOR THE GENERATED TABLES -- do not delete.
%
% The .tex tables written by src/004-analyze-data.R are produced by the
% tinytable package, which builds them with `tabularray` rather than a
% plain tabular environment. Without the block below, \input{} of any
% generated table fails with "Environment talltblr undefined."
%
% This is copied verbatim from tinytable's documentation. If you switch
% the table generator, you can drop it.
% -----------------------------------------------------------------------------
\usepackage{tabularray}
\usepackage[normalem]{ulem}
\UseTblrLibrary{booktabs}
\UseTblrLibrary{siunitx}
\newcommand{\tinytableTabularrayUnderline}[1]{\underline{#1}}
\newcommand{\tinytableTabularrayStrikeout}[1]{\sout{#1}}
\NewTableCommand{\tinytableDefineColor}[3]{\definecolor{#1}{#2}{#3}}

\graphicspath{{../output/}}

\title{Replicating \citet{beaver1968}:\\
       The Information Content of Annual Earnings Announcements}
\author{Your Name\\ACCT 932}
\date{\today}

\begin{document}

\maketitle

\begin{abstract}
\noindent
\guidance{One paragraph. What did you do, and what did you find? Write
this last.}
Replace this with a short abstract.
\end{abstract}


% =============================================================================
\section{Introduction}
% =============================================================================

\guidance{What question did Beaver ask, why did it matter in 1968, and
what does your replication do? Two or three paragraphs is plenty.}

Replace this text.


% =============================================================================
\section{Research design}
% =============================================================================

\guidance{Describe the design in prose: the sample, the event window, the
two outcome measures, and how you handle announcements that fall on
non-trading days.}

\guidance{Be explicit about where you depart from Beaver. The departures
are the interesting part, and the ``DESIGN CHOICE'' comments in
src/002-transform-data.R flag each one.}

\guidance{The design implemented here follows \citet{gowding2024},
Chapter 12 -- cite it when you describe the event-window construction.
The chapter is free online at
\url{https://iangow.github.io/far_book/beaver68.html} and is worth
reading before you write this section.}

Replace this text.

% NOTE: each generated .tex file is a COMPLETE, self-contained table --
% it brings its own \begin{table}, \centering, and \caption. So you
% \input{} it bare. Wrapping it in another \begin{table} produces the
% error "LaTeX Error: Not in outer par mode."
\input{../output/sample-selection.tex}


% =============================================================================
\section{Results}\label{sec:results}
% =============================================================================

\subsection{Descriptive statistics}

\input{../output/descriptives.tex}

\guidance{Interpret the table in a paragraph.}

Replace this text with a paragraph interpreting the table.

\subsection{Replication of Beaver's figures}

\begin{figure}[H]
\centering
\includegraphics[width=0.85\textwidth]{fig1-volume.pdf}
\caption{Trading volume around annual earnings announcements. Compare to
         Figure 1 of \citet{beaver1968}.}
\label{fig:volume}
\end{figure}

\begin{figure}[H]
\centering
\includegraphics[width=0.85\textwidth]{fig2-return-variability.pdf}
\caption{Return variability around annual earnings announcements. Compare
         to Figure 6 of \citet{beaver1968}.}
\label{fig:variability}
\end{figure}

Replace this text.

\subsection{Formal statistical tests}

\guidance{Beaver assessed significance informally. This is the modern
version -- explain what the specification does and why the fixed effects
and the clustered standard errors are there.}

\input{../output/main-results-returns.tex}

\input{../output/main-results-turnover.tex}

Replace this text.

\subsection{Has the effect changed over time?}

\begin{figure}[H]
\centering
\includegraphics[width=0.85\textwidth]{fig3-turnover-by-decade.pdf}
\caption{Median share turnover around annual earnings announcements, by decade.}
\label{fig:turnover}
\end{figure}

\input{../output/by-decade.tex}

Replace this text.


% =============================================================================
\section{Extension: your own partition}
% =============================================================================

\guidance{This section is the heart of the assignment. Everything above
reproduces a result that is already known; this is where you ask a
question of your own. Weight your effort accordingly -- this section
should be the longest in the document.}

\guidance{The by-decade split above is your worked example: it replicates
Figure 1 of \citet{dechow2014}. Here you do the same thing with a
partitioning variable you choose yourself. Two constraints from the
syllabus: firm size is off limits, and your variable must be approved
before you build on it. README.md, under ``Splitting the sample by your
own variable'', lists the four places in the pipeline you need to
change.}

\subsection{The variable, and why you chose it}

\guidance{Name the variable and define it precisely -- the exact
construction, not just the concept. Then give the economic reason to
expect the announcement reaction to differ across it. What is the
mechanism? Cite prior work where it exists.}

Replace this text.

\subsection{Hypothesis}

\guidance{State a directional prediction, and state it before you look at
your figure. Which group should show the larger announcement-window
reaction, and why? Then say what pattern would count as evidence against
you -- a prediction that no result could contradict is not a
hypothesis.}

Replace this text.

\subsection{What you found}

\begin{figure}[H]
\centering
\includegraphics[width=0.85\textwidth]{fig6-my-partition.pdf}
\caption{Announcement-period reaction, partitioned by your own variable.}
\label{fig:mypartition}
\end{figure}

\input{../output/my-partition.tex}

\guidance{Read your own figure honestly. Does it match the prediction you
just made? Be specific about magnitude, not only direction: ``larger'' is
not a finding, ``roughly twice the baseline elevation, against 1.4 times
for the other group'' is. If the prediction failed, say so plainly and
work out why. A sharp prediction that fails is worth more than a vague
one that could not.}

Replace this text.

\subsection{The obvious objection}

\guidance{One paragraph. What is the most serious alternative explanation
for what you found, and what test would separate it from yours? You are
not required to run that test -- you are required to know what it is.}

Replace this text.


% =============================================================================
\section{Discussion questions}
% =============================================================================

\guidance{Answer each question in a paragraph or two. These are graded on
the quality of the reasoning, not on length.}

% Numbered by the enumerate environment rather than by hand, so that
% cutting or adding a question renumbers the rest by itself -- in the
% Word version too. Do not type numbers into the \textbf{} titles.
\begin{enumerate}

\item \textbf{Research questions.}
How do the research questions of \citet{beaver1968} and
\citet{ballbrown1968} differ? If there is overlap, does one paper provide
superior evidence, or are they just different?

\medskip\noindent\textit{Your answer:}

\medskip

\item \textbf{Data differences.}
What differences are there in the data (sample period, frequency, number
of firms) used by \citet{beaver1968} relative to \citet{ballbrown1968}?
Why do these differences exist?

\medskip\noindent\textit{Your answer:}

\medskip

\item \textbf{Price versus volume.}
Which is the better variable for Beaver's research question---price or
volume? Is it helpful to have both?

\medskip\noindent\textit{Your answer:}

\medskip

\end{enumerate}


% =============================================================================
\section{Conclusion}
% =============================================================================

Replace this text.


\bibliography{references}

\end{document}
